\chapter{Architecturen en geheugen}
\label{chap:architecturen}

Volgt later..

% In de voorgaande hoofdstukken is kort al gesproken over een bekende computerarchitectuur, namelijk de Von Neumann-architectuur. Er zijn nog een tweetal andere architecturen die we in dit hoofdstuk gaan behandelen, die te maken hebben met de wijze waarop geheugen wordt behandeld. Wanneer we ons op het gebied van geheugen begeven, halen we ook het concept van virtueel geheugen aan en behandelen we de wijze waarop dit geïmplementeerd wordt.

% \subsection*{Leerdoelen}
% Na het bestuderen van dit hoofdstuk wordt verwacht dat je:
% \begin{itemize}
%     \setlength\itemsep{-1pt}
%     \item[--] de Von Neumann-architectuur kent en begrijpt.
%     \item[--] de Harvard-architectuur kent en begrijpt.
%     \item[--] weet wat virtueel geheugen is.
%     \item[--] weet wat een page is als geheugenconcept.
%     \item[--] weet wat een Memory Management Unit (MMU) is.
%     \newline
% \end{itemize}

% \subsection*{Relevant materiaal}
% Bij dit hoofdstuk hoort hoofdstuk 7 "Processor and Memory Architectures" van het tekstboek van \textbf{Ledin} (\href{https://learning.oreilly.com/library/view/modern-computer-architecture/9781803234519/Text/Chapter\_7.xhtml\#:-:text=7,Processor\%20and\%20Memory\%20Architectures}{link}).


% \section{Computerarchitecturen}
% \studeeraanwijzing{Bestudeer in \textbf{Ledin} vanaf "The von Neumann, Harvard, and modified Harvard architectures" tot "Physical and virtual memory" (\href{https://learning.oreilly.com/library/view/modern-computer-architecture/9781803234519/Text/Chapter\_7.xhtml\#:-:text=The\%20von\%20Neumann,In}{link}).}

% In Ledin laat de auteur doorschemeren dat de \intextnote{Harvard-architectuur} misschien een betere optie zou zijn geweest in vergelijking met de veelgebruikte Von Neumann-architectuur, omdat de eerstgenoemde veiliger zou zijn. Het is waar dat het de Harvard-architectuur immuun is voor het soort code-injectie technieken waar de Von Neumann-architectuur wel vatbaar voor is, maar dat betekent niet dat het hele probleem daarmee is opgelost. Er blijft helaas nog steeds een punt bestaan waar we code moeten vertrouwen voordat we deze inladen om uit te voeren. Voor wat betreft de Harvard-architectuur, zou dit een mogelijk punt van aanval kunnen zijn. Het probleem wordt dus deels naar voren verschoven.

% De volgende link is erg interessant om te weten te komen hoe een bufferoverflow werkt: 
% \url{https://www.youtube.com/watch?v=1S0aBV-Waeo}. Dit maakt overigens geen deel uit van de stof die getoetst zal worden.

% Een versie van de Harvard-architectuur kan ook gevonden worden in bepaalde microcontrollers. De architectuur past hierbij omdat microcontrollers en RAM en \intextnote{ROM} kunnen hebben. ROM staat voor Read Only Memory, en dit betekent dus dat de code wel uit een speciale regio van het geheugen komt waar geen data kan staan.

% \begin{opdracht}
% Maak opdracht 1 van hoofdstuk 7 van \textbf{Ledin} (\href{https://learning.oreilly.com/library/view/modern-computer-architecture/9781803234519/Text/Chapter\_7.xhtml\#:-:text=Exercises}{link}).
% \end{opdracht}

% \section{Geheugen}
% \studeeraanwijzing{Bestudeer in \textbf{Ledin} vanaf "Physical and virtual memory" tot "Paged virtual memory" (\href{https://learning.oreilly.com/library/view/modern-computer-architecture/9781803234519/Text/Chapter\_7.xhtml\#:-:text=Physical\%20and\%20virtual\%20memory,Memory\%20devices\%20in\%20computers}{link}).}

% We slaan het onderdeel over Windows NT paged memory over.

% \studeeraanwijzing{Bestudeer in Ledin vanaf "Memory pools" tot "Summary" (\href{https://learning.oreilly.com/library/view/modern-computer-architecture/9781803234519/Text/Chapter\_7.xhtml\#:-:text=Memory\%20pools,Windows\%20NT\%20categorizes}{link}).}

% \begin{opdracht}
% Maak opdracht 2 van hoofdstuk 7 van \textbf{Ledin} (\href{https://learning.oreilly.com/library/view/modern-computer-architecture/9781803234519/Text/Chapter\_7.xhtml\#:-:text=Exercises}{link}).
% \end{opdracht}
